\newcounter{hyp}
\renewcommand{\thehyp}{H\arabic{hyp}}

\section*{Refined Research Hypotheses}

\noindent\textbf{\refstepcounter{hyp}\thehyp\label{hyp:h1} (Data Structure and Cache Locality).} 
Order books that use index-based layouts (like arrays) will be significantly faster 
than search trees or sorted vectors. This advantage is primarily driven by 
improved CPU cache behaviour during high-frequency updates.

\medskip
\noindent\textbf{Acceptance criteria:} The array implementation shows the lowest mean 
latency in direct mode across most scenarios, and performance counters show 
lower memory-access pressure compared to others.

\bigskip
\noindent\textbf{\refstepcounter{hyp}\thehyp\label{hyp:h2} (Memory Allocation Limits).} 
While memory pooling can reduce specific allocation costs, it will not overcome 
the fundamental performance bottlenecks of a slow data structure.

\medskip
\noindent\textbf{Acceptance criteria:} The pool implementation shows 
statistically significant gains ($p < 0.0125$ with Bonferroni correction) 
over the standard map, but does not reach the performance levels of the 
array where sequential searching is the main bottleneck.

\bigskip
\noindent\textbf{\refstepcounter{hyp}\thehyp\label{hyp:h3} (System-Level Masking).} 
In a networked (gateway) environment, the overhead of the network and system 
queues will account for the vast majority of end-to-end latency, masking the 
algorithmic efficiency of the order book implementation.

\medskip
\noindent\textbf{Acceptance criteria:} For the primary non-vector 
implementations (array, map, pool), the combined transport overhead 
(network and queue) accounts for the significant majority ($>85\%$) 
of total end-to-end latency.
